\chapter*{Introduction}

Les plantes occupent une place essentielle dans la vie quotidienne de l'être humain, que ce soit pour l'alimentation, la santé, l'aménagement des espaces verts ou la préservation de l'environnement. Les pépinières jouent ainsi un rôle fondamental en fournissant des plantes adaptées à différents usages et en contribuant à la biodiversité ainsi qu'à l'équilibre écologique. Cependant, la gestion et l'entretien des plantes nécessitent des connaissances spécifiques, notamment en ce qui concerne leurs besoins, leur croissance et les problèmes pouvant affecter leur santé.  

Avec l'augmentation de l'intérêt pour le jardinage, l'agriculture urbaine et les espaces verts, de nombreux utilisateurs se trouvent confrontés à des difficultés liées au choix des plantes, à leur entretien ou à l'identification des maladies et des carences. Il devient alors nécessaire de proposer des solutions permettant d'accompagner les utilisateurs tout au long du cycle de vie des plantes, depuis leur acquisition jusqu'à leur entretien.  

Dans ce contexte, les pépinières doivent non seulement assurer la disponibilité et la diversité des plantes, mais aussi offrir un accompagnement permettant aux utilisateurs de mieux comprendre et gérer leurs végétaux. L'association de la vente de plantes avec des outils d'aide à l'entretien et au diagnostic constitue une réponse pertinente à ces besoins, en facilitant l'accès aux plantes tout en apportant un soutien adapté. 

 C'est dans cette perspective que s'inscrit notre projet de fin d'année, dont l'objectif est de concevoir une plateforme de pépinière en ligne offrant à la fois la vente de plantes et un accompagnement pour le diagnostic et la recommandation. Ce projet vise à améliorer l'expérience des utilisateurs, à réduire les erreurs d'entretien et à favoriser une meilleure prise en charge des plantes. 

 Le premier chapitre décrit le cadre général du projet. Il présente le contexte de travail, l'étude et la critique de l'existant, les objectifs du projet ainsi que le travail demandé et la méthodologie du travail adoptée.
 Dans le deuxieme chapitre nous abordons l'analyse et spécification des besoins. Il introduit la description du système, l'identification des acteurs 
 